\documentclass{article}
\usepackage{graphicx}
\usepackage{amsmath}
\usepackage{amsfonts}
\usepackage{bookmark}
\usepackage{booktabs}
\usepackage{exercises}
\graphicspath{ {./images/} }
\usepackage{indentfirst}
\usepackage{amsthm}
\usepackage{yhmath}
\usepackage{placeins}
\newcommand*{\oldneg}{\mathord{\sim}}
\newcommand\showdiv[1]{\overline{\smash{)}#1}}
\newcommand\Mydiv[1]{%
\smash{\raise.3ex\hbox{$\big)$}}\mkern-8mu
        \overline{\enspace\strut#1}}
\setlength{\tabcolsep}{5pt}
\theoremstyle{definition}
\newtheorem{definition}{Definition}

\setlength{\parindent}{.5in}
\setlength{\parskip}{12pt}
\renewcommand{\arraystretch}{2}
\newcommand{\?}{\stackrel{?}{=}}

\title{Quiz Answers (pg. 288)}
\begin{document}
\maketitle

\section*{1}
(For 1-9, the book asked to check for extraneous solutions. However, none of these problems have more than one solutions. Thus, do we 
really need to check for extraneous solutions?)\newline
\[ 2^{x+1}=16^{x+2} \] \[ 2^{x+1}=(2^4)^{x+2} \] \[ 2^{x+1}=2^{4x+8} \] \[ x+1=4x+8 \] \[ -3x=7 \] \[ x=-\frac{7}{3} \]

\section*{2}
\[ e^{-x}=4 \] \[ \ln (e^{-x})=\ln 4 \] \[ -x\approx 1.386 \] \[ x=-1.386 \]

\section*{3}
\[ 3^{2x}+5=13 \] \[ 3^{2x}=8 \] \[ \log_3 3^{2x}=\log_3 8 \] \[ 2x=\dfrac{\log 8}{\log 3} \]
\[ 2x\approx 1.89 \] \[ x=0.945 \] 

\section*{4}
\[ 3^{x+1}-5=10 \] \[ 3^{x+1}=15 \] \[ \log_3 \left(3^{x+1}\right)=\log_3 15 \] \[ x+1=\dfrac{\log 15}{\log 3} \]
\[ x\approx 1.46 \]

\section*{5}
\[ \log_4\left( 4x+7\right)=\log_4 11x \] \[ 4^{\log_4\left(4x+7\right)}=4^{\log_4 11x} \]
\[ 4x+7=11x \] \[ 7x=7 \] \[ x=1 \]

\section*{6}
\[ \ln\left(3x-2\right)=\ln 6x \] \[ 3x-2=6x \] \[ x=-\frac{2}{3} \]
However, this answer can be rejected, as \[ \ln\left(6\times-\frac{2}{3}\right)=\ln -4 \] is undefined. 

\section*{7}
\[ \log_3 x=-1 \] \[ 3^{\log_3 x}=3^{-1} \] \[ x=\frac{1}{3} \]

\section*{8}
\[ 6\ln x=30 \] \[ \ln x=5 \] \[ e^{\ln x}=e^5 \] \[ x\approx 148.4 \]

\section*{9}
\[ \log_2\left(x+4\right)=5 \] \[ 2^{\log_2\left(x+4\right)}=2^5 \] \[ x+4=32 \] \[ x=28 \]

\section*{11}
We know that: \[ 4=ab^1 \] and \[ 32=ab^2 \] meaning \[ a=\frac{4}{b} \]
We can then solve for $b$: \[ 32=\frac{4}{b}b^2 \] \[ 32=4b \] \[ b=8 \] and \[ a=\frac{4}{8}=\frac{1}{2} \]
meaning the exponential function in this situation is:
\[ y=\frac{1}{2}(8)^x \]

\section*{12}
We know that: \[ 15=ab^2 \] and \[ 45=ab^3 \] meaning \[ a=\frac{15}{b^2} \]
We can then solve for $b$: \[ 45=\frac{15}{b^2}b^3 \] \[ 45=15b \] \[ b=3 \] and \[ a=\frac{15}{3^2}=\frac{5}{3} \]
meaning the exponential function in this situation is:
\[ y=\frac{5}{3}(3)^x \]

\section*{14}
We know that: \[ 12=a(3)^b \] and \[ 36=a(10)^b \] meaning \[ a=\frac{12}{3^b} \]
We can then solve for $b$: \[ 36=\frac{12}{3^b}(10)^b \] \[ 36=12\times\left(\frac{10}{3}\right)^b \]
\[ 3=(\frac{10}{3})^b \] \[ \log_{\frac{10}{3}} 3=b \] \[ \dfrac{\log 3}{\log\frac{10}{3}}=b \]
\[ b\approx 0.91 \] and \[ a=\frac{12}{3^{0.91}}\approx 4.42 \]
meaning the power function in this situation is about:
\[ y=4.42x^{0.91} \]

\section*{15}
We know that: \[ 4=a(5)^b \] and \[ 51=a(11)^b \] meaning \[ a=\frac{4}{5^b} \]
We can then solve for $b$: \[ 51=\frac{4}{5^b}11^b \] \[ 51=4\left(\frac{11}{5}\right)^b \]
\[ 12.75=\left(\frac{11}{5}\right)^b \] \[ \log_{\frac{11}{5}} 12.75=b \]
\[ \dfrac{\log 12.75}{\log\frac{11}{5}}=b \] \[ b\approx 3.23 \] and \[ a=\frac{4}{5^{3.23}}\approx 0.022 \]
meaning the power function in this situation is:
\[ y=0.022x^{3.23} \]

\section*{16}
We are given that the average weight of the cod can be given by the model:
\[ y=0.51(1.46)^x \] where $y$ is the weight and $x$ is the age in years. Suppose $y=15$. Thus:
\[ 15=0.51(1.46)^x \] \[ \frac{15}{0.51}=1.46^x \] \[ \log_{1.46}\frac{15}{0.51}=x \]
\[ \dfrac{\log\frac{15}{0.51}}{\log 1.46}=x \] \[ x\approx 8.94 \]
meaning the age is about $8.94$ years.

\end{document}
